\documentclass{jsarticle}
\usepackage[dvipdfmx]{graphicx}
\usepackage{ascmac} % itemboxを利用するために必要
\usepackage{url}

%図にEPSファイルを採用した場合，PDFファイルが生成されないことがあるので，PNGやJPEG，PDFファイルを図に採用したほうがよい．

\title{ソフトウェア設計及び実験\\前期・個人開発進捗報告書}

\author{学生番号：6125010806　氏名：清水 統貴}
\date{2026年7月15日}

\begin{document}
\maketitle

\section{開発期間2025年7月9日～7月15日の進捗状況}
%

\subsection{前回からの課題と進展事項}
%
現在までにキャラクターやギミックの画像の実装，新たなギミックの実装，ステージの作成など
を行った．詳細は表\ref{table1}にを参照．

\begin{table}[htb]
\begin{center}
\caption{課題と進展}
\label{table1}
\begin{tabular}{|p{7cm}|p{7cm}|}\hline
課題 & 進展 \\ \hline \hline
キャラクター，ギミックの画像を追加 & 今まで簡易的にしていたキャラクターやギミックに画像を割り当てた．また，背景の画像も変更した． \\ \hline
ブーストの仕様変更 & 以前はブーストボタンを押しながらもゲージを貯めることができるようにしていたが，ブースト中はゲージがたまらないように変更．また，キーボード操作のときはwキーを長押しではなく連打しないとゲージがたまらないように変更．さらに，ゲージを少し溜めてすぐにブーストすることで，擬似的に常時無敵になることができていたので，ゲージが少なすぎるときには，ブーストボタンを押しても，ゲージは減るが，ブースト状態にはならないようにした． \\ \hline
敵キャラクターの実装 & プレイヤーが一定距離近づくとプレイヤーを追尾してくる敵キャラクターを実装した． \\ \hline
ステージ作成 & ブーストを使わなければ攻略できないようにギミックを配置した．しかし，まだあまり面白みがないので，改善の余地がある． \\ \hline
\end{tabular}
\end{center}
\end{table}

\begin{figure}[htb]
\centering
\label{figure1}
\includegraphics[scale=0.3]{window.png}
\caption{ゲームを実行したときの様子}
\end{figure}

\section{発生した課題と解決策}
今回の実装過程において，不具合及びエラーが発生したが，原因を特定しすべて解決した．
不具合と解決策は表\ref{table2}を参照．
\begin{table}[htb]
\begin{center}
\caption{課題と進展}
\label{table2}
\begin{tabular}{|p{7cm}|p{7cm}|}\hline
不具合 & 解決策 \\ \hline \hline
wキー長押しによるゲージの増加 & キーが押されているかどうかだけを判定していたため，1フレーム前のキーの状態を記憶する処理を追加し，現在キーが押されていて，かつ前のフレームでキーが押されていないときだけゲージを増加させるようにした． \\ \hline
キャラクターとゲージの描画ミス & キャラクターやUIの描画の後に背景を描画していたため，描画の順番を見直し，正しい順番で描画するようにした． \\ \hline
UIの視認性低下 & スコア，タイマー，ゲージが背景の色とかぶり視認性が悪かったので，色を変更し，視認性を向上させた． \\ \hline
\end{tabular}
\end{center}
\end{table}


\section{開発期間2026年7月16日～21日の予定}

\subsection{完了させる事項}
ステージに配置するギミックをもう少し面白くしたい．

\subsection{新しく取り組む事項}
タイトルやクリア後の画面などを実装したい．また，時間があれば，新たなギミックやBGM，Joy-Conの振動機能
などもつけてみたい．


\begin{thebibliography}{9}
\bibitem{gemini} Google : Gemini \url{https://gemini.google.com/?hl=ja} (2026年7月15日閲覧)
\end{thebibliography}
\end{document}

